\section{Trace eight pixels through all six chunk types}
Use the eight-by-one header shown earlier. The payload below is sixteen bytes, followed
by the eight-byte end marker. Offsets are measured from the start of the file. All
bytes and slot numbers in the byte columns are hexadecimal; RGBA components are decimal.

\tagpdfsetup{table/header-rows={1}}
\begin{center}
  \small
  \begin{tabular}{rllll}
    Offset & Bytes                   & Command & Reconstructed RGBA & Cache slot  \\
    \hline
    14     & \texttt{c0}             & RUN one & \((0,0,0,255)\)    & \texttt{35} \\
    15     & \texttt{fe ff 00 00}    & RGB     & \((255,0,0,255)\)  & \texttt{32} \\
    19     & \texttt{c0}             & RUN one & \((255,0,0,255)\)  & \texttt{32} \\
    20     & \texttt{35}             & INDEX   & \((0,0,0,255)\)    & \texttt{35} \\
    21     & \texttt{7f}             & DIFF    & \((1,1,1,255)\)    & \texttt{04} \\
    22     & \texttt{a5 97}          & LUMA    & \((7,6,5,255)\)    & \texttt{0b} \\
    24     & \texttt{ff 07 06 05 80} & RGBA    & \((7,6,5,128)\)    & \texttt{16} \\
    29     & \texttt{04}             & INDEX   & \((1,1,1,255)\)    & \texttt{04} \\
  \end{tabular}
\end{center}

The first run installs opaque black in slot fifty-three. That makes the later
\texttt{35} an opaque-black fetch rather than a fetch of the cache's initial transparent
black. The RGBA command changes alpha to 128 without changing RGB. The final index
restores the alpha of its cached \((1,1,1,255)\) pixel; it does not retain 128 from
the immediately preceding pixel. These two transitions test which state each command
actually supplies.

The next input offset is thirty after all eight pixels. The marker occupies offsets
thirty through thirty-seven, so the complete file length is thirty-eight. The final
word \texttt{0xff010101} is stored as bytes \texttt{01 01 01 ff}. Inspecting only
its low three bytes would miss an alpha-restoration error.

\begin{MARSExercise}{Predict a state before running}
  Pause after the LUMA command. State the next input offset, completed pixel count,
  previous RGBA value and the cache entry changed by that command. Then predict the
  same four quantities after the RGBA command. Explain why the final INDEX changes
  alpha even though it has no explicit alpha byte.
\end{MARSExercise}
